This chapter describes the wearable \emph{wristband} ground unit: voice recognition, ESP32
control firmware, OLED feedback, and the ELRS transmitter interface.
The block diagram in Figure~\ref{fig:block-diagram} applies to this subsystem and to its
connection to the aircraft.

%------------------------------------------------------------------------------
\subsection{System Overview}
\label{subsec:wristband-overview}
%------------------------------------------------------------------------------

The wristband system introduces voice recognition as an additional human--machine interface on
top of conventional RC operation.
After the user speaks a command, the VC-02 module performs offline keyword spotting and sends a
one-byte parameter to the ESP32 over UART.
The ESP32 maps that parameter to CRSF channel values, assembles CRSF frames, and forwards them
to the ELRS transmitter module for wireless transmission to the aircraft receiver and flight
controller.
Compared with a handheld transmitter alone, the wearable flow supports offline recognition (no
cloud), targets low link latency through ELRS, and keeps both hands free during demonstrations.

%------------------------------------------------------------------------------
\subsection{ESP32 Main Controller Design}
\label{subsec:voice}
%------------------------------------------------------------------------------

The ESP32 is the core controller: it parses UART tokens from the VC-02, manages flight-state logic,
generates CRSF payloads, and refreshes the OLED.

The firmware uses two hardware UART peripherals with fixed pin assignments (Arduino core
\texttt{HardwareSerial} instances UART1 and UART2):
\begin{itemize}
    \item \textbf{VC-02 link:} RX on GPIO16, TX on GPIO17, 115,200~bps (connected to the module's
          UART1 in wiring tables; maps to \texttt{HardwareSerial(2)} in the sketch).
    \item \textbf{ELRS / CRSF link:} TX on GPIO25, RX on GPIO26, 115,200~bps (\texttt{HardwareSerial(1)}).
\end{itemize}

After each recognized command, the ESP32 updates sixteen CRSF channels (Roll, Pitch, Throttle,
Yaw, arm, mode, and reserved channels) and transmits a complete CRSF RC frame to the ELRS module
every 20~ms so the airborne receiver never sees a link timeout.

%------------------------------------------------------------------------------
\subsection{VC-02 Voice Recognition Module Design}
%------------------------------------------------------------------------------

An AI-Thinker VC-02 offline voice module performs keyword spotting rather than full-sentence ASR\@.
Audio is digitized on-module, filtered (noise suppression, AGC, endpoint detection), and represented
with Mel-frequency cepstral coefficients before template matching against seven stored commands.
When confidence exceeds an internal threshold, the module emits a UART byte to the ESP32.

Table~\ref{tab:vc-map} matches the spoken words to UART payload values used in firmware.
UART format is 8N1 at 115,200~bps.

\begin{table}[htb!]
    \centering
    \caption{Voice command to UART parameter mapping (VC-02 $\rightarrow$ ESP32).}
    \label{tab:vc-map}
    \begin{tabularx}{\linewidth}{l c X}
        \hline
        \textbf{Voice command} & \textbf{UART byte} & \textbf{Firmware intent} \\
        \hline
        takeoff & 1 & Begin takeoff state machine \\
        land    & 2 & Begin landing sequence \\
        forward & 3 & Timed forward pitch pulse \\
        backward & 4 & Timed backward pitch pulse \\
        left    & 5 & Yaw left (hold) \\
        right   & 6 & Yaw right (hold) \\
        stop    & 7 & Brake / hover depending on prior motion \\
        \hline
    \end{tabularx}
\end{table}

\begin{table}[htb!]
    \centering
    \caption{VC-02 module to ESP32 connections.}
    \label{tab:vc02-hw}
    \begin{tabularx}{\linewidth}{l l}
        \hline
        \textbf{VC-02 pin} & \textbf{ESP32 pin} \\
        \hline
        UART TX & GPIO16 (ESP32 RX) \\
        UART RX & GPIO17 (ESP32 TX) \\
        GND     & GND \\
        VCC     & 5~V (per module requirement) \\
        \hline
    \end{tabularx}
\end{table}

%------------------------------------------------------------------------------
\subsection{Voice Command and Drone Control Logic}
%------------------------------------------------------------------------------

The ESP32 translates UART tokens into CRSF microsecond channel values.
Takeoff raises throttle in timed stages, then transitions to a programmed hover PWM; land ramps
throttle down before disarming.
Forward and backward commands apply pitch offsets for fixed-duration pulses so the vehicle inches
without requiring a second utterance; yaw-left and yaw-right bias the yaw channel until another
command arrives.
The stop/brake command issues an opposing pitch pulse if the platform was moving forward or
backward, matching the brake helper states in firmware.

Table~\ref{tab:voice-cmds} summarizes the pulse widths used in the demonstrator binary; mid-stick
neutral is 1500~\textmu s with arming on auxiliary channel~5.

\begin{table}[htb!]
    \centering
    \caption{Voice command set and CRSF channel actions (implementation values).}
    \label{tab:voice-cmds}
    \begin{tabularx}{\linewidth}{c l X}
        \hline
        \textbf{Byte} & \textbf{Command} & \textbf{Action} \\
        \hline
        1 & Takeoff & Arm; THR 1250~\textmu s (1.5~s), 1180~\textmu s (1.5~s), hover 1100~\textmu s \\
        2 & Land    & Ramp THR toward 1050~\textmu s, disarm \\
        3 & Forward & Pitch 1550~\textmu s, elevated THR for 1200~ms \\
        4 & Backward & Pitch 1450~\textmu s, elevated THR for 1200~ms \\
        5 & Left    & Yaw 1450~\textmu s (hold) \\
        6 & Right   & Yaw 1550~\textmu s (hold) \\
        7 & Stop    & Brake pulse or hover \\
        \hline
    \end{tabularx}
\end{table}

Directional commands are ignored until the takeoff routine arms the model, preventing motion
commands on the bench.

%------------------------------------------------------------------------------
\subsection{OLED Display Module Design}
%------------------------------------------------------------------------------

A 128$\times$64 SSD1306 OLED shows action name, armed/flight flags, decoded barometric altitude,
and CRSF status strings (\texttt{Waiting CRSF\ldots}, \texttt{CRSF OK}, etc.).
The module uses I\textsuperscript{2}C (\texttt{Wire}) at address 0x3C; on the ESP32 Arduino core the default
SDA/SCL pins are GPIO21 and GPIO22 unless remapped, matching the wrist strap wiring table used in
the design review.

%------------------------------------------------------------------------------
\subsection{ELRS Transmitter Module Design}
\label{subsec:elrs}
%------------------------------------------------------------------------------

\subsubsection{Design Procedure}

Three 2.4 GHz RC link protocols were evaluated: Futaba S-FHSS, FrSky ACCESS, and
ExpressLRS (ELRS).
Table~\ref{tab:rf-comparison} summarizes the key characteristics of each.

\begin{table}[htb!]
    \centering
    \caption{Comparison of evaluated 2.4 GHz RC link protocols.}
    \label{tab:rf-comparison}
    \begin{tabularx}{\linewidth}{l c c c c}
        \hline
        \textbf{Protocol} & \textbf{Max update rate} & \textbf{Min latency} &
            \textbf{Open source} & \textbf{TX power} \\
        \hline
        Futaba S-FHSS  & 55 Hz    & $\approx$18 ms & No  & Fixed      \\
        FrSky ACCESS   & 150 Hz   & $\approx$7 ms  & No  & Fixed      \\
        ELRS 2.4 GHz   & 1000 Hz  & $<$5 ms        & Yes & Configurable \\
        \hline
    \end{tabularx}
\end{table}

Futaba S-FHSS operates at a fixed 55 Hz and uses a closed firmware ecosystem, precluding
integration with custom transmitter hardware.
FrSky ACCESS supports up to 150 Hz but similarly relies on proprietary firmware.
ELRS is an open-source protocol supporting update rates from 50 Hz to 1000 Hz with a publicly
documented firmware stack that can be compiled and deployed on commodity hardware.
ELRS was selected for its low latency, configurable update rate, and compatibility with custom
transmitter hardware.

For the transmitter unit, two hardware options were considered.
The first is a dedicated ELRS TX module (e.g., RadioMaster Ranger Micro), which provides a
self-contained solution with a preconfigured USB interface.
The second is a RadioMaster RP2 receiver module reflashed with ELRS transmitter firmware.
A dedicated TX module adds unnecessary cost and physical bulk for a fixed ground station
installation.
The RP2 module, when flashed with the \emph{RX as TX} firmware variant, performs identically as a
transmitter while occupying a smaller footprint at lower cost.
The RP2-as-TX approach was adopted.

\subsubsection{Hardware}

The RadioMaster RP2 module contains two primary integrated circuits.
The Espressif ESP8285 is a single-chip microcontroller integrating a 32-bit Tensilica L106 core
and a 2.4 GHz 802.11b/g/n Wi-Fi radio.
In the ELRS firmware context, the ESP8285 executes the protocol stack, manages the CRSF
serial interface to the host device, and controls the RF transceiver over SPI.
The Semtech SX1280 is a 2.4 GHz multi-mode transceiver supporting FLRC (Fast Long-Range
Communication) modulation, which ELRS employs at high update rates to achieve packet air times
below 1 ms and end-to-end round-trip latencies below 5 ms at 250 Hz.

\begin{table}[htb!]
    \centering
    \caption{ELRS transmitter module to ESP32 connections.}
    \label{tab:elrs-hw}
    \begin{tabularx}{\linewidth}{l l}
        \hline
        \textbf{ELRS pin} & \textbf{ESP32 pin} \\
        \hline
        TX (to ESP32 RX) & GPIO26 \\
        RX (from ESP32 TX) & GPIO25 \\
        GND & GND \\
        VCC & 5~V (module requirement) \\
        \hline
    \end{tabularx}
\end{table}

\subsubsection{Firmware Configuration}
\label{subsubsec:fw-config}

Firmware was compiled using the ExpressLRS Configurator application.
Three parameters were fixed at compile time.

\textbf{Firmware target.}
The target
\texttt{Generic 2.4\,GHz /\allowbreak Generic ESP8285 2.4\,GHz RX as TX} was selected.
This variant configures the ELRS firmware to initiate and maintain the RF link, assigning the
module the transmitter role in the ELRS packet timing scheme.
A standard receiver build (\texttt{Generic ESP8285 2.4\,GHz RX}) causes the module to listen
for an external transmitter signal and never initiate packets, making it unsuitable for ground-side
use.

\textbf{Regulatory domain.}
The \texttt{ISM\_2400} domain was selected in preference to \texttt{CE\_LBT}
(Listen-Before-Talk).
The CE\_LBT domain mandates that the transmitter sense channel occupancy before each
transmission and defer if the channel is busy.
At high update rates (500 Hz and 1000 Hz), LBT overhead reduces telemetry throughput; ISM\_2400
avoids that limitation.

\textbf{Binding Phrase.}
A user-defined alphanumeric string is compiled into both the transmitter and receiver firmware
builds for automatic linking at power-on.

\subsubsection{Initial Firmware Flashing}
\label{subsubsec:flashing}

The ESP8285 enters UART boot when GPIO0 is held low at reset.
Typical procedure: hold Boot, apply power, release, then flash at 921,600~bps.
If a USB--UART adapter idles low on TX, power the RP2 before attaching TTL leads so strap pins
sample the idle-high state.
Later updates may use the module WebUI.

\subsubsection{CRSF UART Baud Rate to the ESP32}
\label{subsubsec:crsf}

The ESP32 and RP2 use 115,200~bps 8N1 for CRSF in the wristband firmware build.
A 26-byte frame therefore occupies about 2.26~ms, well within the 20~ms service interval used in
software.

\subsubsection{Telemetry Configuration}
\label{subsubsec:telemetry}

When the RP2 serves as the RF head for the wristband stack, downlink telemetry (battery, barometric
altitude, link statistics) is configured through standard ELRS Lua helpers on full-size radios, or
through CRSF inspection paths during bench bring-up.
For laptop ground station use the WebUI \emph{Use as AirPort Serial device} mode can forward
telemetry at higher USB bridge rates.

\subsubsection{CRSF Framing on the ESP32}
\label{subsubsec:crsf-fw}

Sixteen channel widths in 1000--2000~\textmu s are mapped into CRSF 11-bit words, packed into 22 bytes,
and wrapped with the DVB-S2 CRC in a 26-byte frame addressed to 0xC8 with type~0x16, as required
by the CRSF RC-channels packet definition.
